% !TEX root = ../Main.tex
\chapter{语文 \quad 9:00--11:30}

\section{时间分配表}

\begin{center}
\begin{tabular}{@{} >{\bfseries}c l @{}}
    \toprule
    时间段 & 分块 \\
    \midrule
    9:00--9:25 & 信息类文本阅读 \\
    9:25--9:45 & 文学类文本阅读 \\
    9:45--10:05 & 文言文阅读 \\
    10:05--10:15 & 古代诗歌阅读 \\
    10:15--10:30(40) & 语言文字运用 \\
    10:30(40)--11:30 & 作文 \\
    \bottomrule
\end{tabular}
\end{center}

\section{注意事项}

\rmkb{
    \begin{enumerate}
        \item 语文的时间分布大致的原则为 \textbf{"1 分钟写 1 分"}。在这个原则下可以适当调整各个部分的解题顺序。
        \item 语文的时间分布的\textbf{灵活性主要由作文部分来调整}，宁可挤压作文的时间也不要压缩语言基础题的完整性。无论如何\textbf{不建议先写作文}。
        \item 灵活性也体现在\textbf{语言文字运用}部分。对各部分重点是权衡作答与思考程度。
        \item 答题时，要保证\textbf{各个部分的作答完整性}（每个问题都回答），同时保证\textbf{各个问题的作答完整性}（基本要点全覆盖）。\textbf{第二次作答只会浪费时间}。
    \end{enumerate}
}
